% =====================================================================
%  CHAPTER 9: 卫生服务研究
% =====================================================================
\chapter{卫生服务研究}
\setcounter{chapter}{9}
\chapterintro{
掌握卫生服务研究的概念（\textbf{教学难点}）；掌握卫生服务需要、需求与利用的概念及相互关系（\textbf{教学难点}）；掌握卫生服务的综合评价（\textbf{教学难点}）；熟悉卫生资源、卫生费用的概念及评价指标；了解卫生服务研究的内容、方法和我国进展。
}

% ----- 9.1 概述 -----
\section{卫生服务研究概述}

\begin{defbox}
\textbf{卫生服务研究（Health Services Research）}：\imp{（教学难点*）}是从社会学、管理学、经济学等多学科视角，系统研究\keyword{卫生服务的组织、提供、利用、质量和效果}的科学。其核心是研究如何合理组织和有效利用卫生资源，最大程度地满足人群的卫生服务需要，提高卫生服务的质量和效率。
\end{defbox}

\subsection{卫生服务研究的意义与目的}

\begin{enumerate}[leftmargin=*]
    \item \imp{为卫生政策制定提供科学依据}——了解人群的卫生服务需要和需求、利用模式
    \item \imp{优化卫生资源配置}——发现资源配置不合理之处，提出调整方案
    \item \imp{提高卫生服务质量和效率}——评价卫生服务的质量和效果
    \item \imp{促进卫生服务公平性}——识别服务可及性和利用的差异，减少健康不公平
    \item \imp{评价卫生改革的效果}——对卫生政策、项目和改革的实施效果进行科学评估
\end{enumerate}

\subsection{我国卫生服务研究的进展}

\begin{itemize}[leftmargin=*]
    \item 1981年：中美合作在上海县开展卫生服务研究，开创了我国系统化卫生服务研究先河
    \item 1985年后：逐步在全国推广卫生服务研究方法
    \item 1993年起：每5年开展一次\imp{国家卫生服务调查}（National Health Services Survey），已进行6次（1993、1998、2003、2008、2013、2018年）
    \item "国家卫生服务调查"已成为我国卫生政策制定的重要数据来源
\end{itemize}

\subsection{卫生服务研究的内容}

\begin{tcolorbox}[colback=light, colframe=main, boxrule=0.5pt, arc=4pt, title=\textbf{卫生服务研究的六大内容}, breakable]
\begin{enumerate}[leftmargin=*]
    \item \imp{卫生服务需要与需求}——人群健康状况及主观医疗服务需要
    \item \imp{卫生服务利用}——门诊、住院、预防保健等服务的实际使用
    \item \imp{卫生资源}——人力、物力、财力、信息等资源的配置和分布
    \item \imp{卫生服务评价}——服务效果、效率、质量、公平性和可及性评价
    \item \imp{卫生系统绩效}——卫生系统的整体效率和目标达成情况
    \item \imp{卫生政策与改革}——政策制定、实施和效果评估
\end{enumerate}
\end{tcolorbox}

\subsection{卫生服务研究的方法}

\begin{enumerate}[leftmargin=*]
    \item 常规资料分析——利用常规统计报表和行政记录数据
    \item 家庭卫生服务调查——入户问卷调查获取人群层面的数据
    \item 机构调查——对医疗卫生机构的资源、服务量、费用等进行调查
    \item 定性研究——深入访谈和焦点组讨论了解利益相关者的经验和观点
    \item 综合评价方法——运用统计学和管理学方法进行综合评价
\end{enumerate}

% ----- 9.2 卫生服务需要、需求与利用 -----
\section{卫生服务需要、需求与利用}

\subsection{基本概念（教学难点*）}

\begin{keybox}
\imp{（教学难点*）}\textbf{三个核心概念的区别与联系：}

\begin{defbox}
\textbf{卫生服务需要（Health Need）}：从健康（卫生）角度出发，由医学专业人员根据健康标准判断的\keyword{实际健康需要}。包括：个人察觉到的需要（主观需要）和卫生专业人员判定的需要（客观需要）。

\textbf{卫生服务需求（Health Demand）}：从经济学角度，在\keyword{一定的价格水平下，消费者愿意并且有能力购买}的卫生服务量。需求 = 需要 + 购买意愿 + 购买能力。

\textbf{卫生服务利用（Health Service Utilization）}：人群\keyword{实际使用}卫生服务的情况，包括门诊服务利用、住院服务利用和预防保健服务利用等。
\end{defbox}

\textbf{需要、需求、利用三者的关系：}
\begin{enumerate}[leftmargin=*]
    \item \imp{需要≠需求}：有需要的人不一定有需求（如因经济困难、交通不便未就医）；有需求的人不一定有需要（如不必要的体检和药品需求——"诱导需求"）
    \item \imp{需求≠利用}：有需求的人可能因各种障碍（等待时间过长、交通不便等）未能实际利用服务
    \item \imp{需要→需求→利用}是医疗服务转化的理想路径，但转化过程中存在多种障碍
    \item \imp{未满足的需要（Unmet Need）}：有卫生服务需要但没有转化为需求和利用，是需要重点关注的问题
\end{enumerate}
\end{keybox}

\subsection{影响卫生服务需要与利用的因素}

\begin{enumerate}[leftmargin=*]
    \item \imp{人口学特征}：年龄（老年人和儿童需要更多服务）、性别（女性利用率较高）
    \item \imp{社会经济因素}：收入水平（低收入人群未满足需要更多）、教育程度、职业、医疗保险覆盖
    \item \imp{健康状况}：患病率和疾病严重程度是决定需要的最基本因素
    \item \imp{卫生服务供给}：服务的地理可及性、服务质量、价格和等待时间
    \item \imp{文化和信念}：对健康和疾病的文化认知、就医观念和习惯
    \item \imp{政策与制度}：医疗保险制度、转诊制度、支付方式等政策因素
\end{enumerate}

\subsection{常用指标}

\begin{table}[htbp]
\centering
\caption{卫生服务需要与利用常用指标}
\begin{tabularx}{\textwidth}{lX}
\hline
\textbf{指标类型} & \textbf{具体指标} \\
\hline
\imp{卫生服务需要指标} &
两周患病率 = 两周内患病人数/调查人数 $\times$ 100\%；慢性病患病率；健康自评状况；疾病频率指标 \\
\imp{门诊服务利用指标} &
两周就诊率 = 两周内就诊人数/调查人数 $\times$ 100\%；年人均就诊次数；未就诊率及未就诊原因分析 \\
\imp{住院服务利用指标} &
住院率 = 一年内住院人数/调查人数 $\times$ 100\%；人均住院天数；应住院未住院率及原因 \\
\imp{预防保健利用指标} &
计划免疫覆盖率、产前检查率、住院分娩率、健康体检率等 \\
\hline
\end{tabularx}
\end{table}

% ----- 9.3 卫生资源 -----
\section{卫生资源}

\begin{defbox}
\textbf{卫生资源（Health Resources）}：社会投入到卫生服务中的人力、物力、财力、技术和信息等资源的统称。
\end{defbox}

\subsection{卫生人力资源}

\begin{defbox}
\textbf{卫生人力资源（Health Human Resources）}：接受过卫生专业教育和培训，从事卫生服务相关工作的人员的总称。是卫生资源中\keyword{最核心、最活跃的要素}。

\textbf{研究内容：}卫生人力的数量、质量、结构和分布（城乡、地区、机构层级之间）。
\end{defbox}

\subsection{卫生费用}

\begin{defbox}
\textbf{卫生总费用（Total Health Expenditure, THE）}：一个国家或地区在一定时期内（通常为一年）全社会用于卫生保健支出的货币总和。

\begin{itemize}[leftmargin=*]
    \item \imp{来源}：①政府卫生支出（各级政府财政卫生拨款）；②社会卫生支出（社会医疗保险、商业健康保险、企业办医支出等）；③个人现金卫生支出（居民个人直接支付的医药费用）
    \item \imp{分类}：经常性卫生费用（人员经费、药品、材料等日常费用）和基本建设费用（房屋、设备等投资性费用）
    \item \imp{评价指标}：卫生总费用占GDP的百分比；人均卫生总费用；政府卫生支出占财政支出的百分比；个人现金卫生支出占THE比重（目标为降至30\%以下，降低"灾难性卫生支出"风险）
\end{itemize}
\end{defbox}

% ----- 9.4 卫生服务的综合评价 -----
\section{卫生服务的综合评价（教学难点*）}

\begin{keybox}
\imp{（教学难点*）}\textbf{卫生服务的综合评价模式}是将卫生服务需要、卫生资源投入和卫生服务利用三者结合起来进行综合分析的方法。

\textbf{综合评价矩阵：}根据\imp{需要量（高/低）}与\imp{资源投入（高/低）}与\imp{服务利用（高/低）}三个维度的组合，将不同地区分为八种类型：

\centering
\textbf{表9-2 卫生服务综合评价矩阵}
\bigskip

\begin{tabularx}{\textwidth}{cccX}
\hline
\textbf{需要} & \textbf{资源} & \textbf{利用} & \textbf{综合评价} \\
\hline
高 & 高 & 高 & \textbf{A型——资源充足、利用良好、供需平衡} \\
高 & 高 & 低 & \textbf{B型——资源充足但利用不足，需提高利用效率} \\
高 & 低 & 高 & \textbf{C型——资源不足但利用高，需增加资源投入} \\
高 & 低 & 低 & \textbf{D型——资源不足利用差，需全面加强} \\
低 & 高 & 高 & \textbf{E型——可能存在资源过剩或诱导需求} \\
低 & 高 & 低 & \textbf{F型——资源过剩但利用不足，效率低下} \\
低 & 低 & 高 & \textbf{G型——资源低但利用相对好，效率尚可} \\
低 & 低 & 低 & \textbf{H型——需要、资源、利用均低，视情况调整} \\
\hline
\end{tabularx}
\bigskip

\textbf{综合评价的意义：}
\begin{enumerate}[leftmargin=*]
    \item 帮助发现不同地区卫生服务的主要问题和瓶颈
    \item 为卫生资源配置和服务规划提供科学依据
    \item 综合评价三个维度，避免单一指标的片面评价
    \item 适用于地区间比较和不同时期的纵向比较
\end{enumerate}
\end{keybox}

% ----- 复习题 -----
\section{复习题}

\subsection*{一、选择题}
\begin{enumerate}[leftmargin=*, itemsep=3pt]
    \item 以下哪个概念是指"消费者愿意并有能力购买的卫生服务量"？
    \begin{enumerate}
        \item[(A)] 卫生服务需要
        \item[(B)] 卫生服务需求
        \item[(C)] 卫生服务利用
        \item[(D)] 卫生服务供给
    \end{enumerate}
    \item 我国国家卫生服务调查首次开展于：
    \begin{enumerate}
        \item[(A)] 1981年
        \item[(B)] 1985年
        \item[(C)] 1993年
        \item[(D)] 1998年
    \end{enumerate}
    \item 卫生总费用（THE）的来源不包括：
    \begin{enumerate}
        \item[(A)] 政府卫生支出
        \item[(B)] 社会卫生支出
        \item[(C)] 国际援助
        \item[(D)] 个人现金卫生支出
    \end{enumerate}
    \item 综合评价矩阵中，资源充足但利用不足的类型是：
    \begin{enumerate}
        \item[(A)] A型
        \item[(B)] B型
        \item[(C)] C型
        \item[(D)] D型
    \end{enumerate}
\end{enumerate}

\subsection*{二、名词解释}
\begin{enumerate}[leftmargin=*, itemsep=3pt]
    \item \textbf{卫生服务需要（Health Need）}
    \item \textbf{卫生服务需求（Health Demand）}
    \item \textbf{卫生服务利用（Health Service Utilization）}
    \item \textbf{卫生总费用（THE）}
\end{enumerate}

\subsection*{三、简答题}
\begin{enumerate}[leftmargin=*, itemsep=3pt]
    \item 试述卫生服务需要、需求与利用三者之间的区别与联系。
    \item 简述影响卫生服务需要与利用的主要因素。
    \item 简述卫生服务的综合评价模式。
    \item 简述我国国家卫生服务调查的发展历程。
\end{enumerate}

\subsection*{参考答案要点}

\subsubsection*{一、选择题}
1. (B)；2. (C)；3. (C)；4. (B)

\subsubsection*{二、名词解释}
\begin{enumerate}[leftmargin=*]
    \item \textbf{卫生服务需要}：从健康角度由专业人员判断的实际健康需要，包括主观需要和客观需要。
    \item \textbf{卫生服务需求}：在一定的价格水平下，消费者愿意并有能力购买的卫生服务量。需求 = 需要 + 购买意愿 + 购买能力。
    \item \textbf{卫生服务利用}：人群实际使用卫生服务的情况，包括门诊、住院和预防保健服务利用。
    \item \textbf{卫生总费用（THE）}：一个国家或地区在一定时期内全社会用于卫生保健支出的货币总和，来源包括政府、社会和个人三方面。
\end{enumerate}

\subsubsection*{三、简答题}
\begin{enumerate}[leftmargin=*]
    \item 三者关系：需要≠需求≠利用。需要是健康角度的客观判定；需求需要购买意愿和能力；利用是实际使用。存在的差距：有需要无需求（经济障碍）、有需求无利用（等待时间等障碍）、无需要有需求（诱导需求）。
    \item 六大因素：人口学特征（年龄、性别）、社会经济因素（收入、教育、医保）、健康状况（患病率）、服务供给（可及性、质量、价格）、文化和信念、政策制度（医保、转诊、支付方式）。
    \item 综合评价模式根据\imp{需要量（高/低）}、\imp{资源投入（高/低）}和\imp{服务利用（高/低）}三个维度分为八种类型（A-H型），综合分析供需平衡状况，为资源配置和服务规划提供依据。
    \item 1981年中美合作上海县研究开创先河；1993年起每5年开展一次国家卫生服务调查，至今已6次；成为卫生政策制定的重要数据来源。
\end{enumerate}

\newpage
